\section{Conclusione}
I risultati ottenuti indicano che il modello proposto è in grado di affrontare efficacemente il compito assegnato, pur lasciando spazio a ulteriori miglioramenti. In prospettiva, si suggerisce di:
\begin{itemize}
  \item Aumentare quantità e diversità dei dati di addestramento, così da migliorare la capacità di generalizzazione della rete.
  
  \item Sperimentare architetture neurali alternative, al fine di individuare quella più adatta al riconoscimento dei segnali stradali.
  
  \item Rafforzare l'uso di tecniche di regolarizzazione per ridurre l'\textit{overfitting} e limitare l'eccessivo adattamento ai dati di addestramento.
  
  \item Ottimizzare gli \textit{iperparametri} (ad esempio tasso di apprendimento e dimensione del batch) per massimizzare le prestazioni del modello.
  
  \item Valutare la quantizzazione dei pesi per ridurre la complessità del modello mantenendo un livello di accuratezza adeguato.
  
  \item Analizzare le attivazioni nei diversi strati della rete per approfondire la comprensione del processo di apprendimento.
\end{itemize}

In sintesi, il lavoro svolto costituisce una base solida per sviluppi successivi nel riconoscimento dei segnali stradali mediante modelli di \textit{machine learning}. L'iterazione sperimentale e l'esplorazione di nuove soluzioni restano elementi centrali per affrontare scenari sempre più complessi. Una possibile evoluzione applicativa consiste nel passaggio dalla classificazione per macrocategorie all'identificazione del singolo segnale.
